\makeabstract
    {
    Hot Flashes, Pain, and Body Compassion in Postmenopausal Adults With Multiple Sclerosis
    }
    {
    Isha Patel\textsuperscript{1}, Erin Mistretta\textsuperscript{2}, Jennifer Altman\textsuperscript{2}, Lindsey Knowles\textsuperscript{2}, Sarah Simmons\textsuperscript{2}, Eliza Koch\textsuperscript{2}, Dawn Ehde\textsuperscript{2}
    }
    {
    \textsuperscript{1} Amherst College, Amherst, MA\\
\textsuperscript{2} Department of Rehabilitation Medicine, University of Washington School of Medicine, WA
    }
    {
    Women with multiple sclerosis (MS) across stages of menopause may experience pain related to hormonal changes and psychosocial factors. The current study aims to (1) determine whether postmenopausal women with MS who experience hot flashes differ in pain severity and interference, (2) examine the relationship between hot flash bothersomeness and (a) pain interference and (b) pain intensity among participants who report experiencing hot flashes, and (3) examine whether body compassion moderates the relationship between hot flash bothersomeness and pain outcomes among postmenopausal women with MS.
    }
    {
    Secondary analysis of cross-sectional self-report data collected from menopausal adults with multiple sclerosis (MS) (n= 198). Statistical methods included independent-samples t-test, Pearson product-moment correlations, and multiple linear regressions.
    }
    {
    Women who endorsed hot flashes reported significantly higher pain intensity (M = 59.21, SD = 9.32 vs. M = 53.93, SD = 11.04, p \ensuremath{<} .001, 97.5\% CI [1.99, 8.56]) and interference (M = 58.80, SD = 9.82 vs. M = 54.91, SD = 10.23, p = .008, 97.5\% CI [0.63, 7.16]) than women who did not. Among women who endorsed hot flashes, bothersomeness was positively correlated with pain interference, r(79) = .268, p = .016, 97.5\% CI [.02, .48]. The correlation with pain intensity was similar in magnitude, r(79) = .233, p = .036, 97.5\% CI [-.02, .46] but not statistically significant. Greater hot flash bothersomeness and lower body compassion were independently associated with worse pain outcomes; however, body compassion did not moderate these associations. 
    }
    {
    These findings underscore the importance of considering menopausal symptoms, particularly hot flashes, during clinical assessment of postmenopausal adults with MS. Longitudinal and intervention studies are needed to clarify the directionality and clinical significance of these associations.
    }

    \newpage
    